\documentclass[12pt,a4paper]{article}
\usepackage{amsmath,amssymb,amsthm}
\begin{document}
\section{Names and accents}
Mathematical history is full of names that are hard to typeset. We list a few: Erd\H{o}s, Bol\'yai, Gau\ss, \O rsted, Ma\v{c}ek, G\"odel, Poincar\'e, Cauchy, \'Emile Borel, Fr\'echet, \c{S}en, \r{A}ngstr\"om, Ko\l{}mogorov, Lebesgue, Hadamard, Bi\^anchi, Hurwitz, and \u{S}evchenko.

The \ae sthetic of a proof is a matter of taste; compare \"uber and na\"ive.  Some more marks: \`a la carte, se\~nor, \=a, \.z, Dvo\v{r}\'ak, \k{a}.

\section{Dashes and quotes}
Pages 10--20 cover the case $n\ge2$---the interesting one. The statement is ``every bounded sequence has a convergent subsequence,'' due to Bolzano--Weierstrass. A single quote: `this' is used rarely.

Ties keep things together: see Section~1, Theorem~3, and Mr.~Brown.

\section{Line control}
This line is broken by hand\newline and this one \linebreak[4]continues after a suggested break. Do not break\nolinebreak[4] here.
\end{document}
